\section{Requirements, System Model, and Design Principles}
\label{sec:requirements}

\subsection{Terminology and entities}

The following distinctions are used throughout the report.

\begin{description}[style=nextline]
  \item[Module type and module instance]
  A module type is reusable semantic and mechanical data in a \robotpack{}. A module
  instance is one placed occurrence in a runtime scene.

  \item[Connector type and connector instance]
  A connector type defines compatibility, gender, capture tolerance, orientation,
  physical constraint intent, and policy. A connector instance attaches that meaning to
  a link and local pose on a module type.

  \item[Physical constraint and logical connection]
  A physical constraint is a backend object such as a \mujoco{} weld. A logical
  connection is canonical \modsim{} state created only after the physical constraint
  has been accepted.

  \item[Assembly]
  An assembly is a connected component of the committed module graph. A disconnected
  module is an assembly of size one.

  \item[Canonical state and model view]
  \code{WorldState} is mutable runtime authority. A model view is an immutable,
  JSON-safe projection generated from that authority for algorithms or interfaces.

  \item[Reconfiguration plan and physical route]
  A plan names connector removals and additions. A route describes how a mobile module
  or assembly attempts to realize one plan action through physics.
\end{description}

\subsection{Functional requirements}

\begin{longtable}{L{0.10\textwidth}L{0.31\textwidth}L{0.47\textwidth}}
\caption{Primary requirements and their implemented mechanisms.}
\label{tab:requirements}\\
\toprule
ID & Requirement & Implemented mechanism \\
\midrule
\endfirsthead
\toprule
ID & Requirement & Implemented mechanism \\
\midrule
\endhead
R1 & Portable platform package & Strict split-YAML Robot Pack with relative assets,
stable identifiers, metadata, mappings, and recipes. \\
R2 & Early and profile-aware validation & Authoring and simulation-readiness profiles
with structured issues and cross-reference checks. \\
R3 & Backend independence & Lazy adapter registry, capability negotiation, and no
mandatory simulator dependency in core. \\
R4 & Deterministic docking & Sorted candidates, structured acceptance, guards,
explicit commands, stable connection identifiers, and deterministic events. \\
R5 & Physical/logical consistency & Two-phase commit: backend creates the constraint
before \code{DockCommitted} mutates canonical state. \\
R6 & Dynamic topology & Event-driven connection changes, assembly recomputation,
revision counters, and generated topology multigraphs. \\
R7 & Actuated dynamics & Backend-neutral atomic joint commands with declaration,
capability, finiteness, and limit validation; MuJoCo effort implementation. \\
R8 & Inspectable execution & Immutable Runtime Inspector frames, contiguous event
deltas, graph/status/metrics UI, and native physics viewer from one session. \\
R9 & Safe authoring & Local assets, path containment, strict schemas, transactional
updates, canonical writing, and non-overwriting export. \\
R10 & Extensibility & Registries and typed protocols for backends and model-view
builders without executable code embedded in Robot Packs. \\
\bottomrule
\end{longtable}

\subsection{Non-functional requirements}

\paragraph{Determinism.}
Stable IDs, sorted iteration, canonical YAML, explicit event sequence numbers, and
deterministic candidate ordering make failures reproducible. Physical trajectories may
still depend on solver and floating-point behavior, but semantic arbitration should not
depend on hash or insertion order.

\paragraph{Dependency isolation.}
The base package depends only on Pydantic, Rich, ruamel.yaml, and Typer. \mujoco{} and
Studio reside behind independent optional extras. Importing \code{modsim} must not
import Qt, VTK, or \mujoco{}.

\paragraph{Fail-closed behavior.}
Unsupported physical connection types, absent backend capabilities, exhausted weld
pools, unresolved frames, invalid command modes, and bad cross-references should
produce actionable errors rather than a simulated success.

\paragraph{Bounded memory.}
The event log is intentional history, while model-view caches retain only the latest
result per source and recipe. UI publication is throttled and transfers immutable
snapshots rather than sharing backend objects.

\paragraph{Auditability.}
Every committed dock and undock has a canonical event. Runtime logs are local,
per-launch sidecars. The GUI does not silently persist edits: users explicitly Save or
Export a canonical pack.

\subsection{Authority rules}

The most important architectural decisions can be expressed as authority rules.

\begin{enumerate}
  \item A mechanical asset owns detailed geometry and kinematics. Robot Pack YAML
        should not duplicate an entire link tree.
  \item A Robot Pack owns portable modular semantics. A backend model should not be the
        only place a docking face is named or typed.
  \item \code{WorldState} owns live logical state. A graph widget or backend scene must
        not become an independent connection database.
  \item A backend owns integration and the existence of physical constraints. Core
        code never writes directly into \mujoco{} arrays.
  \item Events own committed semantic history. Candidates and transient acceptance
        state may be ephemeral; topology mutations may not.
  \item Model views own presentation-oriented projections only. They can be rebuilt
        from the pack and world and are never edited to control canonical state.
\end{enumerate}

\subsection{Current non-goals}

The first phase intentionally excludes a new physics solver, a CAD system, a general
autonomous planner, a calibrated magnetic connector model, arbitrary plugin discovery,
and a web application. The APIs are designed so these can be added without replacing
the semantic core. This is particularly important for a future high-performance C++
runtime: Python coordinates authoring and research workflows today, while immutable
data contracts and narrow backend protocols avoid making Python object identity part of
the long-term wire model.

